\pdfmapfile{+lm.map}
\pdfmapfile{+symbols.map}
\documentclass[11pt]{article}
\usepackage[a4paper,margin=26mm]{geometry}
\usepackage[T1]{fontenc}
\usepackage{lmodern,amsmath,amssymb,amsthm,booktabs}
\usepackage{xurl}
\usepackage[hidelinks]{hyperref}
\setlength{\emergencystretch}{3em}
\usepackage{microtype}
\newtheorem{theorem}{Theorem}
\newtheorem{lemma}{Lemma}
\newcommand{\Conf}{\operatorname{Conf}}
\newcommand{\Q}{\mathbb{Q}}
\title{Chromatic number does not determine\\ configuration-space Betti numbers}
\author{A counterexample to TLMC conjecture 00000004091}
\date{14 October 2026}
\begin{document}
\maketitle
\begin{abstract}
The polynomial clause of conjecture 00000004091 is false already for one
particle. The graphs $K_2$ and $C_4$ both have chromatic number $2$, whereas
$\Conf_1(K_2)$ and $\Conf_1(C_4)$ have first Betti numbers $0$ and $1$.
We give the complete elementary argument and a self-contained Lean 4.31.0
verification of the color-minimality and rational cellular-homology
certificates. The subdivision-invariance clause is not challenged.
\end{abstract}

\section{Statement and interpretation}
The source conjecture \cite{tlmc} concerns the unordered configuration space
of $n$ pairwise distinct points of a graph. It asserts both subdivision
invariance of the first Betti number and its expression as a polynomial in
the graph's chromatic number, of degree exactly $n$.
There is no restriction $n\geq 2$ in either the English or Chinese statement.
The polynomial assertion is interpreted in its substantive, universal sense:
for each $n$, one polynomial works for all graphs. Allowing a different
polynomial for each individual graph would impose no relation between these
invariants.

Throughout, a finite graph denotes its usual one-dimensional CW realization,
not its discrete vertex set. Its first Betti number is
$\beta_1(G)=\dim_{\Q}H_1(G;\Q)$.
Both graphs below are finite, simple, nonempty, and connected.

\begin{theorem}
There is no function of chromatic number alone that equals
$\beta_1(\Conf_1(G))$ for every finite simple connected graph $G$.
Consequently the polynomial clause of conjecture 00000004091, and hence
its conjunction with subdivision invariance, is false.
\end{theorem}

\section{The one-particle reduction}
For a topological space $X$, the ordered configuration space is
\[
 F_n(X)=\{(x_1,\ldots,x_n)\in X^n:x_i\ne x_j\text{ for }i\ne j\},
 \qquad \Conf_n(X)=F_n(X)/S_n.
\]
This is the standard definition used for graphs \cite[p.~2]{scheirer}.
When $n=1$, the distinctness requirement has no instances, so $F_1(X)=X^1$.
Evaluation $(x)\mapsto x$ and insertion $x\mapsto(x)$ are continuous inverse
maps for the product topology. The group $S_1$ has one element, so every
orbit is a singleton. The quotient map $q:X^1\to X^1/S_1$ is bijective;
it is a homeomorphism because the quotient-topology definition says that
$U$ in the quotient is open exactly when $q^{-1}(U)$ is open. Thus the map
$[(x)]\mapsto x$ is a natural homeomorphism
\[
                         \Conf_1(X)\cong X.
\]
Homeomorphism invariance of homology reduces the theorem to the two graph
calculations below.

\section{Exact graph calculations}
\subsection{Chromatic numbers}
The graph $K_2$ has vertices $0,1$ and its single edge joins them. Coloring
these vertices $0,1$ gives a proper two-coloring. Since the edge endpoints
must have different colors, one color cannot suffice. Hence $\chi(K_2)=2$.

Let $C_4$ have vertices $0,1,2,3$ and edges
$\{0,1\},\{1,2\},\{2,3\},\{3,0\}$.
The parity coloring $i\mapsto i\bmod 2$ is proper. The presence of an edge
again excludes one color. Thus $\chi(C_4)=2$.

\subsection{Rational cellular homology}
Choose one orientation for every edge. The cellular chain spaces are
$C_1(G;\Q)=\Q^{E(G)}$, $C_0(G;\Q)=\Q^{V(G)}$, and $C_2(G;\Q)=0$.
For an edge $e$ directed from $u$ to $v$, the boundary is
$\partial_1e=v-u$. Therefore
\[
 H_1^{\mathrm{cell}}(G;\Q)
 =\ker\partial_1/\operatorname{im}\partial_2
 =\ker\partial_1.
\]
This is the standard graph chain description in Hatcher
\cite[pp.~99--101]{hatcher}. Cellular homology agrees with singular homology
by Theorem~2.35, and the same comparison with arbitrary coefficients is
explained in the section beginning on p.~153 \cite{hatcher}.

For $K_2$, orient its edge $0\to1$. The boundary matrix is
\[
 B_{K_2}=\begin{pmatrix}-1\\1\end{pmatrix}.
\]
A chain with coefficient $a$ has zero boundary precisely when
$(-a,a)=(0,0)$, which forces $a=0$.
Thus the kernel has its empty basis and $\beta_1(K_2)=0$.

For $C_4$, use the coherent orientations
$0\to1\to2\to3\to0$. In this order the boundary matrix is
\[
 B_{C_4}=\begin{pmatrix}
 -1&0&0&1\\
 1&-1&0&0\\
 0&1&-1&0\\
 0&0&1&-1
 \end{pmatrix}.
\]
The kernel equations for $z=(a,b,c,d)$ are
$d-a=0$, $a-b=0$, $b-c=0$, and $c-d=0$.
They are equivalent to $a=b=c=d$.
Hence every cycle has the unique expression
$a(1,1,1,1)$, and the nonzero all-one chain is a basis.
It follows that $\beta_1(C_4)=1$.

\subsection{Contradiction}
If one function $f$ of chromatic number supplied the claimed Betti number,
these two calculations and the one-particle homeomorphism would give
\[
                f(2)=0\qquad\text{and}\qquad f(2)=1,
\]
which is impossible. In particular, no polynomial of degree exactly $1$
can do so for $n=1$. No assertion about degree or polynomial coefficients
is needed for this contradiction. This proves the theorem.

Subdivision invariance itself is compatible with this argument: subdivision
does not change the topological realization of the graph, and the induced
homeomorphism acts on configuration spaces coordinatewise. It is the
chromatic-determination clause that fails.

\section{Lean verification and its scope}
The project imports only \texttt{Std} and targets Lean 4.31.0. Its finite
graph structure stores oriented endpoints and proves that there are no
loops or repeated undirected edges.
The source \texttt{Conjecture4091/Basic.lean} contains:
\begin{itemize}
\item \texttt{boundary}: the actual incoming-minus-outgoing incidence map,
with edge coefficients in Lean's exact rational type \texttt{Rat}.
\item \texttt{IsCycle}: the kernel equations at every vertex.
\item \texttt{HasFirstBetti G d}: an actual $d$-vector basis of that kernel,
expressed by spanning and uniqueness of rational linear-combination
coefficients. Betti values are not assigned as labels.
\item \texttt{segment\_cycle\_iff} and \texttt{square\_cycle\_iff}: complete
characterizations of the two incidence kernels, followed by basis
certificates \texttt{segment\_firstBetti} and \texttt{square\_firstBetti}.
\item \texttt{HasChromaticNumber}: a legal coloring and a proof that every
smaller color count is impossible. Both witnesses are certified at $2$.
\item {\small\texttt{conjecture\_00000004091\_false}}: negation of the necessary
one-particle consequence that an arbitrary function
$\mathbb N\to\mathbb N$ determines the first Betti number from chromatic
number. The contradiction uses the two certificates at the same input $2$.
\end{itemize}
\subsection{Explicit bridge to the topological statement}
For a stored graph $G$, form $|G|$ by taking one point for each stored
vertex and one copy of $[0,1]$ for each stored edge, identifying $0$ with
its source and $1$ with its target. This is the finite one-dimensional CW
realization with exactly the stored vertices and open edges as cells.
The source and target data of \texttt{segment} give $K_2$; those of
\texttt{square} give the four edges of $C_4$ listed above.

Sending a rational function on the stored edges to the corresponding
linear combination of oriented cellular edges is a vector-space
isomorphism $\Q^{\mathrm{Fin}(G.\mathrm{edges})}\to C_1(|G|;\Q)$;
the analogous map identifies stored vertex functions with $C_0$.
Under these maps, the coefficient of a vertex in the cellular boundary
is exactly the sum on incoming edges minus the sum on outgoing edges.
Thus Lean's \texttt{boundary} is the cellular boundary, not an assigned
numerical invariant. Its \texttt{linearCombination} is the ordinary
rational linear combination in these coordinates. Lean's \texttt{Rat}
represents the field of rational numbers exactly.

A witness for \texttt{HasFirstBetti G d} is therefore a genuine basis
of $\ker\partial_1$: every proposed vector lies in the kernel, every
kernel vector is spanned, and the coefficients are unique. Uniqueness
implies linear independence by applying it to the zero vector.
Consequently its $d$ equals $\dim_\Q\ker\partial_1$.
Because $|G|$ has no cells in dimension two, $\operatorname{im}\partial_2=0$.
The cellular-to-singular comparison cited above now identifies this
kernel with $H_1(|G|;\Q)$. Finally, the homeomorphism from Section~2
identifies it with $H_1(\Conf_1(|G|);\Q)$.
For colorings, choosing an orientation does not affect whether the two
endpoints have different colors, so \texttt{HasChromaticNumber} gives the
usual chromatic minimum of the underlying undirected graph.
These identifications verify the semantic link for each witness.

The formal dimension certificates use the standard cellular definition of
first Betti number for a graph. The quotient topology of $\Conf_1$, its
homeomorphism with $G$, and the general cellular-to-singular comparison
are mathematical bridges proved or cited above; they are not machine-checked
by this Std-only project. No such bridge is introduced as a Lean axiom.
The proof is a semantic counterexample in the graph cellular model,
with this topological boundary stated explicitly.

The final theorem and basis theorems use only the standard Lean logical
axioms \texttt{propext}, \texttt{Classical.choice}, and \texttt{Quot.sound}.
There are no \texttt{sorry} placeholders, \texttt{native\_decide}, or custom
axioms. The included build logs record successful compilation and a separate
\texttt{leanchecker} run.

\section{Independent computation and eligibility}
The separate \texttt{empirical.py} script performs exact Gaussian elimination
over Python rational fractions and exhaustive coloring searches. It confirms
the incidence ranks $1$ and $3$ and kernel dimensions $0$ and $1$ of the two
witnesses. It also enumerates all $1{,}099$ labeled simple graphs with one
through five vertices. Among connected bipartite graphs in this range,
it finds first Betti numbers $0,1,2$ at the same chromatic number $2$.
These computations are independent corroboration; they are not premises
of the Lean proof.

At the live check on 14 October 2026 at 19:46 UTC, conjecture 00000004091
had \texttt{proven=false} and \texttt{disproven=false}; its solutions path
was absent. All-state searches for the full ID and for Betti/subdivision
found the earlier, closed, unmerged submission PR~292 \cite{prior}.
The short-ID and configuration/chromatic searches returned no matches;
these inconsistent search results are why the prior PR was also fetched
directly. The response snapshot is retained in
\path{evidence/eligibility-00000004091.json}. This records a check at that
time, not acceptance or priority over earlier or later work.

\section{Earlier submission and the present repair}
The counterexample idea is already present in orionsheep's PR~292 of
3 October 2026: at one particle, a path and an even cycle have the same
chromatic number but different first Betti numbers. We credit that earlier
observation and do not claim it as a new mathematical discovery.
That PR used $P_4$ and $C_4$; the present proof uses the shorter path $K_2$
and the same cycle.

The organizer closed PR~292 without merging it. The recorded reason was
that its Lean statement only conjoined elementary integer equalities and
inequalities, with no graph, configuration-space, or Betti-number
objects connecting the arithmetic to the conjecture \cite{prior-review}.
Inspection of that PR's source confirms the issue:
\texttt{b1\_C4} and \texttt{b1\_P4} are just arithmetic identities,
\texttt{same\_chromatic} is the conjunction $2=2\wedge2=2$, and
\texttt{conjecture\_refuted} merely assembles those statements.
This is a formalization gap; the elementary paper counterexample is sound.

The present project supplies the missing finite simple graphs, proper
colorings and minimality proofs, incidence boundary on rational chains,
and kernel bases with unique coefficients. Its concluding theorem refutes
a universally quantified dependence on chromatic number using those
certificates. The exact bridges from this cellular model to the topological
conjecture are given in Section~4.1. They remain paper mathematics, rather
than Lean theorems. Thus this submission repairs the graph/chromatic/Betti
semantic gap but does not claim a full Lean formalization of general
configuration-space topology. Whether this disclosed boundary meets the
competition's review standard remains for its reviewers to decide.
A detailed comparison is included in \texttt{PRIOR\_SUBMISSION.md}.

\clearpage
\begin{thebibliography}{9}
\bibitem{tlmc}
The Last Math Competition, conjecture 00000004091.
\url{https://github.com/The-Last-Math-Competition/The-Last-Math-Competition/blob/main/conjectures/00000004091.md}.
Source blob: \texttt{73f51cfa010d743d7f2363be23ee139f9c819802}.
\bibitem{hatcher}
Allen Hatcher, \emph{Algebraic Topology}, Cambridge University Press, 2002.
Chapter 2, pp.~99--101; Theorem~2.35, pp.~139--140; homology with
coefficients, p.~153.
\url{https://pi.math.cornell.edu/~hatcher/AT/ATch2.pdf}.
\bibitem{scheirer}
Steven Scheirer, \emph{Topological complexity of unordered configuration
spaces of certain graphs}, 2019, p.~2.
\url{https://arxiv.org/pdf/1905.12838}.
\bibitem{prior}
orionsheep, PR~292, \emph{Disproof of conjecture 00000004091},
3 October 2026, closed without merging.
\url{https://github.com/The-Last-Math-Competition/The-Last-Math-Competition/pull/292}.
\bibitem{prior-review}
lidangzzz, review comment on PR~292, 3 October 2026.
\url{https://github.com/The-Last-Math-Competition/The-Last-Math-Competition/pull/292#issuecomment-5968931449}.
\end{thebibliography}
\end{document}
